%\clearpage
\section{Funciones continuas en espacios métricos compactos; homeomorfismos}

\begin{teorema}
Sea $f:A\to Y$ continua en $A$ y $B\subseteq A$ compacto.
Entonces $f(B)$ es compacto en $Y$.
\end{teorema}

\begin{proof}
Sea $\mathcal{F}=\{U_i\}$ un cubrimiento abierto de $f(B)$.
Como cada $U_i$ es abierto y $f$ es continua, $f^{-1}(U_i)$ es abierto en $A$.

Además, $B\subseteq\bigcup_i f^{-1}(U_i)$, pues si $x\in B$ entonces $f(x)\in f(B)$ y por tanto $f(x)\in U_i$ para algún $i$, lo que implica $x\in f^{-1}(U_i)$.

Así, $\{f^{-1}(U_i)\}$ es un cubrimiento abierto de $B$.
Por compacidad de $B$, existen $U_1,\dots,U_n\in\mathcal{F}$ tales que
\[
B\subseteq\bigcup_{i=1}^n f^{-1}(U_i).
\]
Aplicando $f$, obtenemos
\[
f(B)\subseteq\bigcup_{i=1}^n U_i.
\]
Luego $f(B)$ es compacto.
\end{proof}

\begin{teorema}
Sea $f:A\to\R$ continua y $B\subseteq A$ compacto.
Entonces existen $p,q\in B$ tales que
\[
f(p)=\inf f(B),
\qquad
f(q)=\sup f(B).
\]
\end{teorema}

\begin{remark}
En el teorema anterior, $f(p)$ y $f(q)$ se llaman respectivamente el valor mínimo absoluto y el valor máximo absoluto de $f$ en $B$.
\end{remark}

\begin{teorema}[Teorema del valor extremo]
Sea $f:[a,b]\to\R$ continua en $[a,b]$.
Entonces $f$ alcanza su mínimo y su máximo absolutos en $[a,b]$.
\end{teorema}

\begin{definicion}
Sea $f:A\to\R$.
\begin{enumerate}
    \item $p\in A$ es punto de máximo absoluto si $f(x)\le f(p)$ para todo $x\in A$.
    \item $p\in A$ es punto de mínimo absoluto si $f(p)\le f(x)$ para todo $x\in A$.
\end{enumerate}
\end{definicion}

\begin{example}
La función $f:(0,1)\to\R$, $f(x)=x$, es continua en $(0,1)$, pero no alcanza ni máximo absoluto ni mínimo absoluto.
\end{example}

\begin{teorema}
Sea $f:A\to Y$ inyectiva y continua. Si $A$ es compacto, entonces
\[
f^{-1}:f(A)\to A
\]
es continua.
\end{teorema}

\begin{proof}
Como $f$ es inyectiva, $f^{-1}$ está bien definida sobre $f(A)$.
Sea $C\subseteq A$ cerrado en $A$. Como $A$ es compacto, $C$ es compacto.
Por continuidad de $f$, $f(C)$ es compacto; en particular, es cerrado en $f(A)$.

Ahora,
\[
(f^{-1})^{-1}(C)=f(C),
\]
que es cerrado en $f(A)$. Por lo tanto, $f^{-1}$ es continua.
\end{proof}

\begin{definicion}[Homeomorfismo]
Sea $f:S\to T$ una función entre espacios métricos y supongamos que es inyectiva.
Decimos que $f$ es un homeomorfismo si
\[
f:S\to f(S)
\]
es continua y
\[
f^{-1}:f(S)\to S
\]
es continua.
\end{definicion}

\begin{example}
Los espacios métricos $\R$ y $(-\pi/2,\pi/2)$ (con la métrica de subespacio) son homeomorfos, pues
\[
\tan:(-\pi/2,\pi/2)\to\R
\]
es un homeomorfismo.
\end{example}
